\summarychapter{总结 关于特征为 0 的域上的有限维李代数的某些性质}{总结}{关于特征为 0 的域上的有限维李代数的某些性质.}


设 g 是一个李代数, r 是其根基, n 是其最大幂零理想, s 是其幂零根基, t 是相对于 Killing 型的 g 的正交补. 于是 r, n, s, t 是特征理想, 且 \(r \supset t \supset n \supset s\) .

\begin{enumerate}
\def\labelenumi{(\Roman{enumi})}
\tightlist
\item
  以下性质中的任意一个都刻画半单李代数:
\end{enumerate}

\begin{enumerate}
\def\labelenumi{(\arabic{enumi})}
\tightlist
\item
  \(r = \{0\}\) ; (2) \(n = \{0\}\) ; (3) \(t = \{0\}\) ; (4) g 的每个交换理想都是零; (5) 代数 g 同构于简单李代数的一个乘积; (6) g 的每个有限维表示都是半单.
\end{enumerate}

\begin{enumerate}
\def\labelenumi{(\Roman{enumi})}
\setcounter{enumi}{1}
\tightlist
\item
  以下性质中的任意一个都刻画约化李代数:
\end{enumerate}

\begin{enumerate}
\def\labelenumi{(\arabic{enumi})}
\tightlist
\item
  \(\mathfrak{s} = \{0\}\) ; (2) \(\mathfrak{r}\) 是 \(\mathfrak{g}\) 的中心; (3) \(\mathcal{D}\mathfrak{g}\) 是半单; (4) \(\mathfrak{g}\) 是一个半单代数与一个交换代数的乘积; (5) \(\mathfrak{g}\) 的伴随表示是半单的; (6) \(\mathfrak{g}\) 有一个有限维表示, 使得相应的双线性型非退化; (7) \(\mathfrak{g}\) 有一个有限维半单忠实表示.
\end{enumerate}

\begin{enumerate}
\def\labelenumi{(\Roman{enumi})}
\setcounter{enumi}{2}
\tightlist
\item
  以下性质中的任意一个都刻画可解李代数:
\end{enumerate}

\begin{enumerate}
\def\labelenumi{(\arabic{enumi})}
\tightlist
\item
  \(\mathcal{D}^{p}\mathfrak{g} = \{0\}\) 当 p 充分大时; (2) 存在 g 的理想组成的递减序列 \(g = g_{0} \supset g_{1} \supset \cdots \supset g_{n} = \{0\}\), 使得代数 \(g_{i}/g_{i+1}\) 是交换的; (3) 存在一个递减序列
\end{enumerate}

\[
\mathfrak {g} = \mathfrak {g} _ {0} ^ {\prime} \supset \mathfrak {g} _ {1} ^ {\prime} \supset \dots \supset \mathfrak {g} _ {n ^ {\prime}} ^ {\prime} = \{0 \}
\]

\(\mathfrak{g}\) 的子代数, 使得 \(\mathfrak{g}_{i+1}^{\prime}\) 是 \(\mathfrak{g}_i^{\prime}\) 的理想, 且 \(\mathfrak{g}_i^{\prime}/\mathfrak{g}_{i+1}^{\prime}\) 是交换的; (4) 存在一个递减序列 \(\mathfrak{g} = \mathfrak{g}_0'' \supset \mathfrak{g}_1'' \supset \cdots \supset \mathfrak{g}_{n''}'' = \{0\}\) 的 \(\mathfrak{g}\) 的子代数, 使得 \(\mathfrak{g}_{i+1}''\) 是 \(\mathfrak{g}_i''\) 中余维 1 的理想; (5) \(\mathfrak{g} \supset \mathcal{D}\mathfrak{g}\) ; (6) \(\mathcal{D}\mathfrak{g}\) 是幂零的.

\begin{enumerate}
\def\labelenumi{(\Roman{enumi})}
\setcounter{enumi}{3}
\tightlist
\item
  以下性质中的任意一个都刻画幂零李代数:
\end{enumerate}

\begin{enumerate}
\def\labelenumi{(\arabic{enumi})}
\tightlist
\item
  \(\mathcal{C}_{g}^{p} = \{0\}\) 当 p 充分大时; (2) \(\mathcal{C}_{p}\mathfrak{g} = \mathfrak{g}\) 当 p 充分大时; (3) 存在 g 的理想组成的递减序列 \(g = g_{0} \supset g_{1} \supset \cdots \supset g_{p} = \{0\}\), 使得 \([g, g_{i}] \subset g_{i+1}\) ; (4) 存在 g 的理想组成的递减序列 \(g = g_{0}' \supset g_{1}' \supset \cdots \supset g_{p'}' = \{0\}\), 使得 \([g, g_{i}'] \subset g_{i+1}'\), 且 \(g_{i}'/g_{i+1}'\) 是 1-维的; (5) 存在一个整数 i, 使得
\end{enumerate}

\[
\left(\operatorname{ad} x _ {1}\right) \circ \left(\operatorname{ad} x _ {2}\right) \circ \dots \circ \left(\operatorname{ad} x _ {i}\right) = 0
\]

对所有 \(x_{1},\ldots ,x_{i}\), 它们都在 \(\mathfrak{g}\) 中 ; (6) 对所有 \(x\in \mathfrak{g}\) , ad \(x\) 都是幂零的.

\begin{enumerate}
\def\labelenumi{(\Alph{enumi})}
\setcounter{enumi}{21}
\tightlist
\item
  g 交换 \(\Rightarrow\) g 幂零 \(\Rightarrow\) g 的 Killing 型为零 \(\Rightarrow\) g 可解.
\end{enumerate}

g 交换 \(\Rightarrow\) g 约化.

g 半单 \(\Rightarrow\) g 约化.

\textbf{(VI) r 的刻画:}

\begin{enumerate}
\def\labelenumi{(\arabic{enumi})}
\tightlist
\item
  r 是 g 的最大可解理想; (2) r 是使 g/r 为半单的最小理想; (3) r 是使 g/r 为半单的唯一可解理想; (4) r 是相对于 Killing 型的 \(\mathcal{D}\mathfrak{g}\) 的正交补.
\end{enumerate}

\textbf{(VII) n 的刻画:}

\begin{enumerate}
\def\labelenumi{(\arabic{enumi})}
\tightlist
\item
  n 是 g 的最大幂零理想; (2) n 是 r 的最大幂零理想; (3) n 是 \(x \in r\) 中使 \(\operatorname{ad}_{\mathfrak{g}} x\) 为幂零的元素的集合; (4) n 是 \(x \in r\) 中使 \(\operatorname{ad}_{\mathfrak{r}} x\) 为幂零的元素的集合; (5) n 是 g 的最大理想, 使得对所有 \(x \in n\), \(\operatorname{ad}_{\mathfrak{g}} x\) 都是幂零的; (6) n 是 \(x \in g\) 中使 \(\operatorname{ad}_{\mathfrak{g}} x\) 属于由 1 和 \(\operatorname{ad}_{\mathfrak{g}} y\,(y \in g)\) 生成的结合代数的根基的元素的集合.
\end{enumerate}

\textbf{(VIII) s 的刻画:}

\begin{enumerate}
\def\labelenumi{(\arabic{enumi})}
\tightlist
\item
  \(\mathfrak{s}\) 是 \(\mathfrak{g}\) 的有限维简单表示的核的交; (2) \(\mathfrak{s}\) 是 \(\mathfrak{g}\) 的有限维半单表示的核中最小的一个; (3) \(\mathfrak{s}\) 是 \(\mathfrak{g}\) 的有限维表示的最大幂零性理想的交; (4) \(\mathfrak{s}\) 是 \(\mathfrak{g}\) 的最小理想, 使得 \(\mathfrak{g} / \mathfrak{s}\) 是约化的; (5) \(\mathfrak{s} = \mathfrak{r} \cap \mathcal{D}\mathfrak{g}\) ; (6) \(\mathfrak{s} = [\mathfrak{r}, \mathfrak{g}]\) ; (7) \(\mathfrak{s}\) 是相对于与 \(\mathfrak{g}\) 的有限维表示相关的双线性型的 \(\mathfrak{g}\) 的正交补的交.
\end{enumerate}


